% GENERATED FILE. Edit canonical lessons, then run npm run generate:books.
% canonical-source-hash: fnv1a64:2ea814904d00c822
% canonical-lessons: UR-C423-sauda, UR-C423-bhao, UR-C423-dam, UR-C423-riayat, UR-C423-sel

\chapter{Rates and Deals}
\label{ch:ur-rates-and-deals}
\hlchaptermodality{423}

\begin{chapteropening}
\textbf{By the end of this chapter:} \emph{I can ask the going rate, hear the price and get a discount.}

Complete the last lesson of chapter 423: I can ask the going rate, hear the price and get a discount.
\end{chapteropening}

\section[\texorpdfstring{\ur{سودا} (saudā)}{سودا (saudā)}]{\ur{سودا} (saudā) — a deal; groceries}

\label{lesson:UR-C423-sauda}



\begin{quote}
Before the new one: say the Urdu for to collect (money owed), then the Urdu for to give back.
\end{quote}

\subsection*{You'll want to know: \ur{سودا}}
\textbf{\ur{سودا}} — \emph{saudā} — \textquotedblleft{}a deal; groceries\textquotedblright{}.

\begin{quote}
\ur{بازار} \ur{میں} \ur{ہم} \ur{نے} \ur{سودا} \ur{خریدا۔}
\end{quote}

\begin{quote}
\emph{bāzār meṅ ham ne saudā kharīdā.}
\end{quote}

\begin{quote}
We bought our groceries at the market.
\end{quote}

Masculine.

\begin{cousinweb}
\textbf{\ur{سودا}} (\emph{saudā}) is usually traced to Persian \emph{saudā}, \textquotedblleft{}trade, a bargain\textquotedblright{}; in Urdu it is a deal, and also the day's shopping.
\end{cousinweb}

\begin{grammarlens}[title={What you've built}]
One more word for getting around.
\end{grammarlens}

\subsection*{Your turn}
\begin{itemize}
\raggedright
  \item \emph{Say it:} \emph{saudā}
  \item \emph{Say it:} \emph{saudā}, once more
  \item \emph{Say it:} \emph{bāzār meṅ ham ne saudā kharīdā.} — We bought our groceries at the market.
  \item \emph{Recall:} say the Urdu for to collect (money owed), then the Urdu for to give back, then say \emph{saudā} again
\end{itemize}

\begin{tcolorbox}[breakable,colback=teal!4,colframe=teal!35!black,title={Before you move on}]
What does \ur{سودا} mean? (\textquotedblleft{}A deal; groceries\textquotedblright{}.) Say it once more.
\end{tcolorbox}
\section[\texorpdfstring{\ur{بھاؤ} (bhāo)}{بھاؤ (bhāo)}]{\ur{بھاؤ} (bhāo) — the going rate; a price per unit}

\label{lesson:UR-C423-bhao}



\begin{quote}
Before the new one: say the Urdu for to give back, then the Urdu for a deal.
\end{quote}

\subsection*{You'll want to know: \ur{بھاؤ}}
\textbf{\ur{بھاؤ}} — \emph{bhāo} — \textquotedblleft{}the going rate; a price per unit\textquotedblright{}.

\begin{quote}
\ur{آج} \ur{آلو} \ur{کا} \ur{بھاؤ} \ur{کیا} \ur{ہے؟}
\end{quote}

\begin{quote}
\emph{āj ālū kā bhāo kyā hai?}
\end{quote}

\begin{quote}
What are potatoes going for today?
\end{quote}

Masculine.

\begin{cousinweb}
\textbf{\ur{بھاؤ}} (\emph{bhāo}) is inherited: CDIAL lists the Hindi cousin \emph{bhāu}, \textquotedblleft{}a price\textquotedblright{}, under Sanskrit \emph{bhāvá}, \textquotedblleft{}being, state\textquotedblright{} (CDIAL 9475).
\end{cousinweb}

\begin{grammarlens}[title={What you've built}]
One more word for getting around.
\end{grammarlens}

\subsection*{Your turn}
\begin{itemize}
\raggedright
  \item \emph{Say it:} \emph{bhāo}
  \item \emph{Say it:} \emph{bhāo}, once more
  \item \emph{Say it:} \emph{āj ālū kā bhāo kyā hai?} — What are potatoes going for today?
  \item \emph{Recall:} say the Urdu for to give back, then the Urdu for a deal, then say \emph{bhāo} again
\end{itemize}

\begin{tcolorbox}[breakable,colback=teal!4,colframe=teal!35!black,title={Before you move on}]
What does \ur{بھاؤ} mean? (\textquotedblleft{}The going rate; a price per unit\textquotedblright{}.) Say it once more.
\end{tcolorbox}
\section[\texorpdfstring{\ur{دام} (dām)}{دام (dām)}]{\ur{دام} (dām) — the asking price; a price}

\label{lesson:UR-C423-dam}



\begin{quote}
Before the new one: say the Urdu for a deal, then the Urdu for the going rate.
\end{quote}

\subsection*{You'll want to know: \ur{دام}}
\textbf{\ur{دام}} — \emph{dām} — \textquotedblleft{}the asking price; a price\textquotedblright{}.

\begin{quote}
\ur{دکاندار} \ur{نے} \ur{کم} \ur{دام} \ur{مانگے۔}
\end{quote}

\begin{quote}
\emph{dukāndār ne kam dām māṅge.}
\end{quote}

\begin{quote}
The shopkeeper asked a low price.
\end{quote}

Masculine.

\begin{cousinweb}
\textbf{\ur{دام}} (\emph{dām}) is inherited: CDIAL lists the Hindi cousin \emph{dām}, \textquotedblleft{}a copper coin, money, a price\textquotedblright{}, under Sanskrit \emph{dramma}, \textquotedblleft{}a coin\textquotedblright{}, which was itself taken from Greek \emph{drachmē} (CDIAL 6622).
\end{cousinweb}

\begin{grammarlens}[title={What you've built}]
One more word for getting around.
\end{grammarlens}

\subsection*{Your turn}
\begin{itemize}
\raggedright
  \item \emph{Say it:} \emph{dām}
  \item \emph{Say it:} \emph{dām}, once more
  \item \emph{Say it:} \emph{dukāndār ne kam dām māṅge.} — The shopkeeper asked a low price.
  \item \emph{Recall:} say the Urdu for a deal, then the Urdu for the going rate, then say \emph{dām} again
\end{itemize}

\begin{tcolorbox}[breakable,colback=teal!4,colframe=teal!35!black,title={Before you move on}]
What does \ur{دام} mean? (\textquotedblleft{}The asking price; a price\textquotedblright{}.) Say it once more.
\end{tcolorbox}
\section[\texorpdfstring{\ur{رعایت} (riāyat)}{رعایت (riāyat)}]{\ur{رعایت} (riāyat) — a discount; a concession}

\label{lesson:UR-C423-riayat}



\begin{quote}
Before the new one: say the Urdu for the going rate, then the Urdu for the asking price.
\end{quote}

\subsection*{You'll want to know: \ur{رعایت}}
\textbf{\ur{رعایت}} — \emph{riāyat} — \textquotedblleft{}a discount; a concession\textquotedblright{}.

\begin{quote}
\ur{دکاندار} \ur{نے} \ur{ہمیں} \ur{رعایت} \ur{دی۔}
\end{quote}

\begin{quote}
\emph{dukāndār ne hameṅ riāyat dī.}
\end{quote}

\begin{quote}
The shopkeeper gave us a discount.
\end{quote}

Feminine.

\begin{cousinweb}
\textbf{\ur{رعایت}} (\emph{riāyat}) is usually traced, through Persian, to Arabic \emph{riʿāya}, \textquotedblleft{}regard, consideration\textquotedblright{}; in a shop it is the price taken off.
\end{cousinweb}

\begin{grammarlens}[title={What you've built}]
One more word for getting around.
\end{grammarlens}

\subsection*{Your turn}
\begin{itemize}
\raggedright
  \item \emph{Say it:} \emph{riāyat}
  \item \emph{Say it:} \emph{riāyat}, once more
  \item \emph{Say it:} \emph{dukāndār ne hameṅ riāyat dī.} — The shopkeeper gave us a discount.
  \item \emph{Recall:} say the Urdu for the going rate, then the Urdu for the asking price, then say \emph{riāyat} again
\end{itemize}

\begin{tcolorbox}[breakable,colback=teal!4,colframe=teal!35!black,title={Before you move on}]
What does \ur{رعایت} mean? (\textquotedblleft{}A discount; a concession\textquotedblright{}.) Say it once more.
\end{tcolorbox}
\section[\texorpdfstring{\ur{سیل} (sel)}{سیل (sel)}]{\ur{سیل} (sel) — a sale (at low prices)}

\label{lesson:UR-C423-sel}



\begin{quote}
Before the new one: say the Urdu for the asking price, then the Urdu for a discount.
\end{quote}

\subsection*{You'll want to know: \ur{سیل}}
\textbf{\ur{سیل}} — \emph{sel} — \textquotedblleft{}a sale (at low prices)\textquotedblright{}.

\begin{quote}
\ur{اس} \ur{ہفتے} \ur{کپڑوں} \ur{کی} \ur{سیل} \ur{ہے۔}
\end{quote}

\begin{quote}
\emph{is hafte kapṛoṅ kī sel hai.}
\end{quote}

\begin{quote}
There's a sale on clothes this week.
\end{quote}

Feminine.

\begin{cousinweb}
\textbf{\ur{سیل}} (\emph{sel}) is borrowed from English \emph{sale}.
\end{cousinweb}

\begin{grammarlens}[title={What you've built}]
One more word for getting around.
\end{grammarlens}

\subsection*{Your turn}
\begin{itemize}
\raggedright
  \item \emph{Say it:} \emph{sel}
  \item \emph{Say it:} \emph{sel}, once more
  \item \emph{Say it:} \emph{is hafte kapṛoṅ kī sel hai.} — There's a sale on clothes this week.
  \item \emph{Recall:} say the Urdu for the asking price, then the Urdu for a discount, then say \emph{sel} again
\end{itemize}

\begin{tcolorbox}[breakable,colback=teal!4,colframe=teal!35!black,title={Before you move on}]
What does \ur{سیل} mean? (\textquotedblleft{}A sale (at low prices)\textquotedblright{}.) Say it once more.
\end{tcolorbox}
